% ============================================================
% РАЗДЕЛ «ИСПОЛЬЗОВАНИЕ ПО НАЗНАЧЕНИЮ»
% ============================================================
{\color{unimain}\section[Использование по назначению]{Использование по назначению}\label{sec:using}}

{\color{unimain}\subsection{Эксплуатационные ограничения}}
Климатические условия эксплуатации должны соответствовать требованиям документа \ManualUnitGeneral.

{\color{unimain}\subsection{Подготовка устройства к использованию}}
Меры безопасности при подготовке устройства, порядок проведения внешнего осмотра устройства перед использованием, требования к монтажу устройства, подключение к персональному компьютеру и устройству «ЮНИТ-ИЧМ» приводится в документации \ManualUnitGeneral.

{\color{unimain}\subsection{Работа с устройством}}
\begin{enumerate}
\item Взаимодействие с устройством осуществляется при помощи подключаемого устройства «ЮНИТ-ИЧМ» или через персональный компьютер с установленным программным обеспечением «ЮНИТ-СЕРВИС-РЗА». Руководство по устройству «ЮНИТ-ИЧМ» представлено в документе \ManualUnitHMI. Руководство по использованию программного обеспечения представлено в документе \ManualUnitSRZA.

% \textcolor{red}{Здесь указать какие текущие значения аналоговых величин и где посмотреть. А также наверно либо здесь картинка с Деревом меню либо ссылка на приложение Дерево меню }
\item Описание работы с устройством (включение, способы управления устройством, обновление встроенного программного обеспечения, неисправности и методы их устранения) приводится в документации \ManualUnitGeneral.
\end{enumerate}

% ============================================================
% РАЗДЕЛ «ТЕХНИЧЕСКОЕ ОБСЛУЖИВАНИЕ И РЕМОНТ»
% ============================================================
{\color{unimain}\section[Техническое обслуживание и ремонт]{Техническое обслуживание и ремонт}\label{sec:techrepair}}

{\color{unimain}\subsection{Техническое обслуживание устройства}}
\begin{enumerate}
\item \sloppy Процесс технического обслуживания устройств релейной защиты регламентирован приказом Минэнерго России «Об утверждении Правил технического обслуживания устройств и комплексов релейной защиты и автоматики и внесении изменений в требования к обеспечению надежности электроэнергетических систем, надежности и безопасности объектов электроэнергетики и энергопринимающих установок».
\item Техническое обслуживание устройств «ЮНИТ-М500» должно выполняться в соответствии с документом \ManualUnitMU.
\end{enumerate}

{\color{unimain}\subsection{Текущий ремонт}}
\begin{enumerate}
\item Устройство серии «ЮНИТ-М500», как было описано выше, выполнено с использованием микропроцессорных технологий, поэтому необходимость в его текущем ремонте отсутствует. Работоспособность и безотказность работы устройства обеспечивается функционированием системы самодиагностики.
\item Углубленная самодиагностика аппаратных средств, а также состояние программной памяти, памяти данных и буферной памяти производится при включении устройства. Поэтому при возникновении сообщений о критических или некритических ошибках от системы самодиагностики, даже при возможности устранения их своими силами, следует обратиться на предприятие-изготовитель для консультации или проведения регламентных работ.
\item Поводом для обращения на предприятие-изготовитель являются также признаки нехарактерного нагрева или перегрева отдельных элементов, участков или частей устройства.
\item Вышедшие из строя устройства должны проходить ремонт на предприятии-изготовителе, обеспечивающем гарантийное и послегарантийное обслуживание, адрес которого указан в паспорте устройства. В качестве ЗИП, если это предусмотрено договором поставки, могут предоставляться устройства «ЮНИТ-М500» с необходимыми конфигурациями, позволяющие оперативно произвести замену вышедших из строя устройств.
\item Перечень сигналов самодиагностики приведен в документации \ManualUnitGeneral.
\end{enumerate}

% ============================================================
% РАЗДЕЛ «ТРАНСПОРТИРОВАНИЕ, ХРАНЕНИЕ И КОНСЕРВАЦИЯ»
% ============================================================
{\color{unimain}\section[Транспортирование, хранение и консервация]{Транспортирование, хранение и консервация}\label{sec:trans}}

Более подробно об условиях транспортирования и хранения устройства можно узнать в соответствующем разделе документации \ManualUnitGeneral.

% ============================================================
% РАЗДЕЛ «УТИЛИЗАЦИЯ»
% ============================================================
{\color{unimain}\section[Утилизация]{Утилизация}\label{sec:util}}

\begin{enumerate}[label=\arabic{section}.\arabic*, labelsep=4pt, leftmargin=0pt, itemindent=57pt] % Добавил, чтобы не было 6.0.1
\item Демонтаж и утилизация устройства не требуют применения специальных мер безопасности и выполняются без применения специальных приспособлений и инструментов.
\item Более подробно о способах утилизации устройства можно узнать в соответствующем разделе документации \ManualUnitGeneral.
\end{enumerate}

% ============================================================
% РАЗДЕЛ «ГАРАНТИЯ ИЗГОТОВИТЕЛЯ»
% ============================================================
{\color{unimain}\section[Гарантия изготовителя]{Гарантия изготовителя}\label{sec:garantee}}

\begin{enumerate}[label=\arabic{section}.\arabic*, labelsep=4pt, leftmargin=0pt, itemindent=57pt] % Добавил, чтобы не было 7.0.1
\item Гарантийный срок эксплуатации устройства составляет 36 месяцев со дня ввода устройства в эксплуатацию, но не более 40 месяцев с даты производства предприятием-изготовителем или с момента проследования через государственную границу при поставках на экспорт.
\item Более подробно о гарантийных обязательствах изготовителя устройства можно узнать в соответствующем разделе документации \ManualUnitGeneral.
\end{enumerate}

% ============================================================
% РАЗДЕЛ «ТЕХНИЧЕСКАЯ ПОДДЕРЖКА»
% ============================================================
{\color{unimain}\section[Техническая поддержка]{Техническая поддержка}\label{sec:support}}

Сервисный центр ООО «Юнител Инжиниринг» принимает заявки от потребителей о качестве функционирования оборудования и ПО, оказывает консультационные услуги при поиске и устранении неисправностей.

Заявки от потребителей принимаются по телефонной связи (тел.: +7 (495) 165-99-98) и (или) по электронной почте:
\begin{itemize}
	\item \href{mailto:rza@uni-eng.ru}{\color{unimain}\uline{rza@uni-eng.ru}}
	\item \href{mailto:info@uni-eng.ru}{\color{unimain}\uline{info@uni-eng.ru}}
\end{itemize}

Форму заявки можно скачать по ссылке:
\begin{itemize}
	\item \href{https://uni-eng.ru/support/servisnyy-tsentr/}{\color{unimain}\uline{https://uni-eng.ru/support/servisnyy-tsentr/}}
\end{itemize}